\section{Discussion}

The comparison separates sensitivity from precision: MaskedBiMamba's higher precision and lower false-alarm rate accompany a lower event recall, a more conservative decision at the common threshold. For an alert system, model selection therefore depends on the relative cost of missed falls and false notifications.

Removing bidirectional processing or temporal attention reduces window and video F1 and increases false alarms, which supports their role in interpreting fall motion. Quality information has a less uniform effect: confidence input improves video F1 and reduces false alarms, while the model without quality pooling attains a higher clean video F1 and removing frame masking raises clean window F1 further. The masking ratio and pooling rule therefore require tuning for the desired balance between sensitivity and false positives.

The strongest advantage appears under missing frames: MaskedBiMamba loses much less window F1 than TCNTE as frame loss increases, and the same pattern appears on Le2i. Adjacent poses provide temporal redundancy that the combined architecture can use when some observations disappear. Random joint loss and complete lower-body removal are harder for both models, since these perturbations remove spatial information across the sequence and leave fewer observations from which to infer the fall trajectory. Testing component ablations under the same perturbations can separate the contribution of each mechanism to this robustness difference.

The system performs pose extraction and temporal inference near the camera, with local records supporting disconnected operation and the cloud providing centralized event history. False alarms remain frequent at the primary threshold. Training on difficult non-fall activities and evaluating the alert rules on continuous recordings are priorities for reducing unnecessary notifications under sustained monitoring.

Relative to the initial edge-feature/cloud-classification plan, the final design also places temporal classification on the Raspberry Pi. This extends the original monitoring objective to disconnected operation: hardware acceleration supports near-real-time inference, local storage retains predictions, and the cloud receives the backlog after recovery. The proposed skeleton/geometry fusion developed into a temporal model with pose-quality features, evaluated through component ablations. HTTP and SQLite implement the event exchange and storage originally assigned to MQTT and PostgreSQL. The resulting system meets the local-processing and record-recovery objectives; the evaluation does not establish a reliable alarm rate for continuous care settings.
